\section{Ablation Plan}
\label{sec:ablations}

Each ablation in \cref{tab:ablations} changes one thing against a named reference system and is decided by a metric fixed in advance. All decoder ablations run with the three seeds of \cref{sec:protocol} and are compared as mean and standard deviation; a difference smaller than the spread across seeds is reported as no difference. Every row also says what the outcome that goes against us would mean.

\begin{table*}[t]
\centering
\small
\begin{tabular}{@{}p{0.17\textwidth}p{0.30\textwidth}p{0.45\textwidth}@{}}
\toprule
Ablation & Hypothesis & Expected interpretation (deciding metric in bold) \\
\midrule
Frozen zero-shot against LoRA ($\lambda = 0$) &
H4, H5: task adaptation is needed; zero-shot is far below the encoder on all three scores and LoRA recovers most of the weighted F1 gap. &
\textbf{Weighted F1 and \HardCons{}}, single-pass. A large gain means generic instruction tuning does not carry the four-way directional scheme. If zero-shot came close to LoRA, the adapters would add little and the comparison with DeBERTa would rest on the backbone alone. A high zero-shot \SoftCons{} with low weighted F1 is the always-\EQ{} pattern, not knowledge. \\
\addlinespace
$\lambda = 0$ against $\lambda > 0$ &
H6: the consistency term raises single-pass \SoftCons{} without lowering weighted F1. &
\textbf{Single-pass \SoftCons{}, with weighted F1 as the constraint.} Higher \SoftCons{} at equal weighted F1 supports H6. Higher \SoftCons{} at lower weighted F1 means the term buys agreement with correctness, possibly by moving pairs to \EQ{} or \NO{}, which are their own reverses; the consistent-but-wrong pair count and the predicted-label counts show that. No change means the both-order cross-entropy already teaches the rule. \\
\addlinespace
Single-pass against symmetrized decoding &
H7: averaging the two orders raises \HardCons{} and weighted F1, not only \SoftCons{}. &
\textbf{\HardCons{} and weighted F1.} \SoftCons{} is not evidence here, because symmetrized decoding raises it by design. If only \SoftCons{} moves, the averaging repairs agreement and leaves correctness alone. A drop in weighted F1 would mean the weaker order drags the stronger one down. No training is needed: both decodings come from the same trained adapters. \\
\addlinespace
With and without the reversed copies of the negatives &
The assumption $\Rev(\NO{}) = \NO{}$ helps \NO{} F1, or at least does not hurt it. &
\textbf{\NO{} F1}, plus the number of reversible pairs answered \NO{}. Without the copies, the second cross-entropy term and the consistency term skip the 478 negatives. If \NO{} F1 rises when the copies are removed, the assumption is wrong for this data; the training file can barely test it, with a single real reversed negative (\path{negatives_with_reversed_copy.csv} in \path{experiments/2026-11-14-decoder-preliminary-track1/tables/}). \\
\addlinespace
Option letters against label words &
The result does not depend on how the label is verbalised. &
\textbf{Weighted F1}, zero-shot and LoRA. Same prompt, but the answer is scored on the first token of each label word instead of A to D. A large gap zero-shot that closes after LoRA means the verbaliser matters only before adaptation. A gap that survives LoRA means our conclusions are tied to the letter format. \\
\addlinespace
0.5B against 1.5B backbone &
The gain of the proposed components does not depend on model size. &
\textbf{Difference of differences}: the gain in \HardCons{} and weighted F1 from $\lambda = 0$ to the chosen $\lambda$, at each size. A gain only at 0.5B suggests the components compensate for a weak model; a gain only at 1.5B suggests they need a model that already reads direction. \\
\addlinespace
Symmetrized decoding on DeBERTa's saved probabilities &
The gain of symmetrized decoding is a general property of the reversal rule, not specific to decoders. &
\textbf{\HardCons{} and weighted F1} of DeBERTa, single-pass against symmetrized. A similar gain means the decoding step is model-agnostic and the decoder-specific claim is the training loss only. No gain on the encoder, which already has \SoftCons{} 0.877, would mean averaging helps only a model that disagrees with itself more often. \\
\bottomrule
\end{tabular}
\caption{Ablation plan. Each row is compared against the system it changes, on the development set, with three seeds where training is involved.}
\label{tab:ablations}
\end{table*}

The last row needs no new training. DeBERTa's per-label probabilities were saved for all three seeds, but outside the repository, in Fawwad's \path{artifacts/} folder and on the Modal volume, so Fawwad runs that ablation. Symmetrized decoding applies to the encoder exactly as in \cref{eq:sym}, with its probabilities in place of the decoder's.

Two rows are not ablations of a trained component. The verbaliser row tests the measurement, and the size row tests how far the conclusions generalise. We keep both because a reader cannot judge the main result without them.
